\chapter{Summary and Outlook}

The \ac{XPS} measurements of \ac{QA} on the Ag(100) surface revealed significant differences between the two phases. First, the C1s spectra revealed distinct components for chemically different C atom, which match to the stoichiometry of the \ac{QA} molecule. The same holds true for the N1s and O1s spectra regarding existence of the hydrogen bonds. This allows to assign the different components of the spectra to specific atoms in the molecule for both phases.

Second, the \ac{BE} of all components, except N1, in the $\beta$-phase are significantly lower than in the $\alpha$-phase, which indicates a stronger molecule-substrate interaction in the $\beta$-phase. This is in line with the fact that the $\beta$-phase has a commensurate structure, which is generally related to a stronger molecule-substrate interaction, while the $\alpha$-phase only has a \ac{pol} structure.

Third, the N1s spectra of the $\alpha$-phase show only one component, which indicates that all N atoms have the same chemical environment and thus form hydrogen bonds. In contrast, the N1s spectra of the $\beta$-phase show two different components with an area ratio of 1:1. The N1 component in the $\beta$-phase is at the same \ac{BE} as the N1 component in the $\alpha$-phase, while the N2 component in the $\beta$-phase is shifted by 2.2~\si{eV} to lower \ac{BE}. The N1 component in the $\beta$-phase can thus be assigned to the N atoms that form hydrogen bonds, while the N2 component can be assigned to the N atoms that do not form hydrogen bonds due to the deprotonation of the N atoms. 

The same holds true for the O1s spectra, where the O1 component can be assigned to the O atoms that form hydrogen bonds, while the O2 component can be assigned to the O atoms that do not form hydrogen bonds due to the deprotonation of the N atoms. Consequently, the \ac{XPS} results indicate that there are two hydrogen bonds per molecule in the $\alpha$-phase, but only one hydrogen bond per molecule in the $\beta$-phase, which contradicts the previously proposed 1.5 hydrogen bonds per molecule in the $\beta$-phase.\autocite{Humberg2020,Humberg2025}

Last, the second component of the N1s and O1s spectra are also present in the \ac{XPS} spectra if the sample is exposed to the X-rays for a long time. This indicates that the X-rays cause a deprotonation of the N atoms, which is a known phenomenon in \ac{XPS} measurements. 



Furthermore, the \ac{NIXSW} measurements revealed the adsorption heights of the different atoms in the \ac{QA} molecule for both phases. In the $\alpha$-phase, the adsorption heights of the C and N atoms are similar and approximately 3.0~\si{\angstrom}, while the O atoms are significantly lower at approximately 2.7~\si{\angstrom}. This indicates a stronger molecule-substrate interaction for the O atoms, which may be the reason for the \ac{pol} structure of the $\alpha$-phase. Furthermore, it indicates that the hydrogen bonds are not this relevant for the adsorption height of the O atoms. 

In the $\beta$-phase, the adsorption heights of all atoms (except N1 \& N2) are closer to the surface than in the $\alpha$-phase. The C atoms are approximately at 2.8~\si{\angstrom}, while the O atoms are again located closer to the surface at approximately 2.6~\si{\angstrom}. This is in line with the stronger molecule-substrate interaction in the $\beta$-phase, which is also indicated by the lower \ac{BE} of all components in the \ac{XPS} spectra and the commensurate structure. 

The N atoms show the most significant change in adsorption height between the two phases. The N1 atoms, which form hydrogen bonds, are located at approximately 3.0~\si{\angstrom} in the $\alpha$-phase, and do not change their adsorption height in the $\beta$-phase. This indicates that the N atoms that form hydrogen bonds are not this relevant for the molecule-substrate interaction and their adsorption height is mainly influenced by the hydrogen bonds.

In contrast, the N2 atoms, which do not form hydrogen bonds, are located at approximately 2.4~\si{\angstrom} in the $\beta$-phase. This indicates that the N atoms that do not form hydrogen bonds are the most relevant for the molecule-substrate interaction. This is supposed to be the reason for the commensurate structure of the $\beta$-phase, which is generally related to a strong molecule-substrate interaction. Furthermore, it indicates that the deprotonation of the N atoms is the main reason for the stronger molecule-substrate interaction in the $\beta$-phase, which compensates the energy loss due to the missing hydrogen bonds.


Finally, the results of thxe \ac{XPS} and \ac{NIXSW} measurements indicate that the previously proposed structure models for \ac{QA} on the Ag(100) surface are not correct, especially for the $\beta$-phase. Thus, new structure models for both phases are proposed based on the results of this thesis, which are discussed in detail in \autoref{chap:models}. The underlying unit cell and molecular arrangement of the two phases are still the same, but the deprotonated N atoms are now taken into account, which leads to a lower number of hydrogen bonds in the $\beta$-phase. In addition, the adsorption heights of the different atoms are now included in the structure models.


In conclusion, the results of this thesis provide new insights into the phase transition of \ac{QA} on the Ag(100) surface, especially regarding the role of the hydrogen bonds and the molecule-substrate interaction. However, there are still open questions regarding the exact mechanism of the phase transition and the exact positions of the deprotonated N atoms in the $\beta$-phase.

The first question could be investigated by performing temperature-dependent \ac{STM} measurements, which may show the rearrangement of the \ac{QA} molecules while the phase transition occurs. The second question could be investigated by performing \ac{NIXSW} measurements for different Bragg planes, which allows the triangulation of the e.g. N atoms and thus the determination of their exact adsorption site.


\cleardoublepage